\SourceLeaf{87}lum fere usum habeat, tamen propter simplicitatem atque generalitatem memoratu dignum est.

\emph{Numerus quicunque $A$ per numerum primum $2m+1$ non divisibilis, huius primi residuum est vel non-residuum, prout $A^m\equiv+1$ vel $\equiv-1\,(\operatorname{mod.}2m+1)$.}

Sit enim pro modulo $2m+1$ in systemate quocunque numeri $A$ index $a$, eritque $a$ par quando $A$ est residuum ipsius $2m+1$, impar vero quando $A$ non-residuum. At numeri $A^m$ index erit $ma$, \emph{i. e.} $\equiv0$ vel $\equiv m\,(\operatorname{mod.}2m)$, prout $a$ par vel impar. Hinc denique $A^m$ in priori casu erit $\equiv+1$, in posteriori vero $\equiv-1\,(\operatorname{mod.}2m+1)$. V. artt. 57, 62.

\emph{Ex.}\quad $3$ ipsius $13$ est residuum quia $3^6\equiv1\,(\operatorname{mod.}13)$, $2$ vero ipsius $13$ non-residuum, quoniam $2^6\equiv-1\,(\operatorname{mod.}13)$.

At quoties numeri examinandi mediocriter sunt magni, hoc criterium ob calculi immensitatem prorsus inutile erit.

\begin{center}\emph{Disquisitiones de numeris primis quorum residua aut non-residua sint numeri dati.}\end{center}

\Article{107}
Facillimum quidem est, proposito modulo, omnes assignare numeros, qui ipsius residua sunt vel non-residua. Scilicet si ille numerus ponitur $=m$, determinari debent quadrata, quorum radices semissem ipsius $m$ non superant, sive etiam numeri his quadratis secundum $m$ congrui (ad praxin methodi adhuc expeditiores dantur), tuncque omnes numeri horum alicui secundum $m$ congrui, erunt residua ipsius $m$, omnes autem numeri nulli istorum congrui erunt non-residua. --- At quaestio inversa, \emph{proposito numero aliquo, assignare omnes numeros quorum ille sit residuum vel non-residuum}, multo altioris est indaginis. Hoc itaque problema, a cuius solutione illud quod in art. praec. nobis proposuimus pendet, in sequentibus perscrutabimur, a casibus simplicissimis inchoantes.

\begin{center}\emph{Residuum $-1$.}\end{center}

\Article{108}
\textsc{Theorema.}\quad \emph{Omnium numerorum primorum formae $4n+1$, $-1$ est residuum quadraticum, omnium vero numerorum primorum formae $4n+3$, non-residuum.}

\emph{Ex.}\quad $-1$ est residuum numerorum $5,\allowbreak 13,\allowbreak 17,\allowbreak 29,\allowbreak 37,\allowbreak 41,\allowbreak 53,\allowbreak 61.73,\allowbreak 89,\allowbreak 97$ etc., e quadratis numerorum $2,\allowbreak 5,\allowbreak 4,\allowbreak 12,\allowbreak 6,\allowbreak 9,\allowbreak 23,\allowbreak 11,\allowbreak 27,\allowbreak 34,\allowbreak 22$ etc. respective ori\SourceLeaf{88}undum; contra non-residuum est numerorum $3,\allowbreak 7,\allowbreak 11,\allowbreak 19,\allowbreak 23,\allowbreak 31,\allowbreak 43,\allowbreak 47,\allowbreak 59,\allowbreak 67,\allowbreak 71,\allowbreak 79,\allowbreak 83$ etc.

Mentionem huius theor. iam in art. 64 fecimus. Demonstratio vero facile ex art. 106 petitur. Etenim pro numero primo formae $4n+1$ est $(-1)^{2n}\equiv1$, pro numero autem formae $4n+3$ habetur $(-1)^{2n+1}\equiv-1$. Convenit haec demonstratio cum ea quam l. c. tradidimus. Sed propter theorematis elegantiam atque utilitatem non superfluum erit, alio adhuc modo idem ostendisse.

\Article{109}
Designemus complexum omnium residuorum numeri primi $p$, quae ipso $p$ sunt minora, excluso residuo $0$, per literam $C$, et quoniam horum residuorum multitudo semper $=\frac{p-1}{2}$, manifestum est, eam fore parem, quoties $p$ sit formae $4n+1$, imparem vero, quoties $p$ sit formae $4n+3$. Dicantur, ad instar art. 77, ubi de numeris in genere agebatur, \emph{residua socia} talia, quorum productum $\equiv1$ (mod.$p$); manifesto enim si $r$ est residuum, etiam $\frac{1}{r}$(mod.$p$) residuum erit. Et quoniam idem residuum plura socia inter residua $C$ habere nequit, patet omnia residua $C$ in classes distribui posse, quarum quaevis bina residua socia contineat. Iam perspicuum est, si nullum residuum daretur, quod sibi ipsi esset socium, \emph{i. e.} si quaevis classis bina residua \emph{inaequalia} contineret, omnium residuorum numerum fore duplum numeri omnium classium; quodsi vero aliqua dantur residua sibi ipsis socia \emph{i. e.} aliquae classes quae unicum tantum residuum aut, si quis malit, idem residuum bis continent, posita harum classium multitudine $=a$, reliquarumque multitudine $=b$; erit omnium residuorum $C$ numerus $=a+2b$. Quare quando $p$ est formae $4n+1$, erit $a$ numerus par; quando autem $p$ est formae $4n+3$, erit $a$ impar. At numeri ipso $p$ minores alii, quam $1$ et $p-1$, sibi ipsis socii esse nequeunt (vid. art. 77); priorque $1$ certo inter residua occurrit; unde in priori casu $p-1$ (seu quod hic idem valet, $-1$) debet esse residuum, in posteriori vero non-residuum; alias enim in illo casu foret $a=1$, in hoc autem $=2$, quod fieri nequit.

\Article{110}
Etiam haec demonstratio ill. Eulero debetur, qui et priorem primus invenit V. \emph{Opusc. Anal. T. I. p. 135.} --- Facile quisquis videbit eam similibus principiis innixam esse, ut demonstratio nostra secunda theor. Wilsoniani art. 77. Si \SourceLeaf{89}vero hoc theorema supponere velimus, facilius adhuc demonstratio exhiberi poterit. Scilicet inter numeros $1,2,3\ldots p-1$ erunt $\frac{p-1}{2}$ residua quadratica ipsius $p$ totidemque non-residua; quare non-residuorum multitudo erit par, quando $p$ est formae $4n+1$; impar quando $p$ est formae $4n+3$. Hinc productum ex omnibus numeris $1,2,3\ldots p-1$ in priori casu erit residuum, in posteriori non-residuum (art. 99). At productum hoc semper $\equiv-1$(mod.$p$); adeoque etiam $-1$ in priori casu residuum, in posteriori non-residuum erit.

\Article{111}
Si itaque $r$ est residuum numeri alicuius primi formae $4n+1$, etiam $-r$ huius primi residuum erit, omnia autem talis numeri non-residua, etiam signo contrario sumta non-residua manebunt\HistoricalNote{*)}{Quando igitur de numero quocunque loquemur quatenus numeri formae $4n+1$ residuum vel non-residuum est, ipsius signum omnino negligere sive etiam signum anceps $\pm$ ipsi tribuere poterimus.}. Contrarium evenit pro numeris primis formae $4n+3$, quorum residua quando signum mutatur, non-residua fiunt et vice versa. vid. art. 98.

Ceterum facile ex praecedentibus derivatur regula generalis: $-1$ \emph{est residuum omnium numerorum qui neque per $4$ neque per ullum numerum primum formae $4n+3$ dividi possunt; omnium reliquorum non-residuum.} V. artt. 103 et 105.

\begin{center}\emph{Residua $+2$ et $-2$.}\end{center}

\Article{112}
Progredimur ad residua $+2$ et $-2$.

Si ex tabula II colligimus omnes numeros primos quorum residuum est $+2$, hos habebimus: $7,17,23,31,41,47,71,73,79.89,97$. Facile autem animadvertitur, inter hos numeros nullos inveniri formarum $8n+3$ et $8n+5$.

Videamus itaque, num haec inductio ad certitudinem evehi possit.

Primum observamus quemvis numerum compositum formae $8n+3$ vel $8n+5$ necessario factorem primum alterutrius formae $8n+3$ vel $8n+5$, involvere; manifesto enim e solis numeris primis formarum $8n+1$, $8n+7$, alii numeri quam qui sunt formae $8n+1$ vel $8n+7$, componi nequeunt. Quodsi itaque inductio nostra generaliter est vera, nullus omnino numerus formae \SourceLeaf{90}$8n+3$, $8n+5$ dabitur, cuius residuum $+2$; sicque nullus certe numerus huius formae infra $100$ exstat, cuius residuum sit $+2$. Si autem ultra hunc limitem tales numeri reperirentur, ponamus minimum omnium $=t$. Erit itaque $t$ vel formae $8n+3$ vel $8n+5$; $+2$ ipsius residuum erit, omnium autem numerorum similium minorum non-residuum. Ponatur $2\equiv aa$(mod.$t$) poteritque $a$ ita semper accipi ut sit impar simulque $<t$, (habebit enim $a$ ad minimum duos valores positivos ipso $t$ minores quorum summa $=t$, quorumque adeo alter par alter impar v. artt. 104. 105). Quo facto sit $aa=2+tu$, sive $tu=aa-2$, eritque $aa$ formae $8n+1$, $tu$ igitur formae $8n-1$, adeoque $u$ formae $8n+3$ vel $8n+5$, prout $t$ est formae posterioris vel prioris. At ex aequatione $aa=2+tu$ sequitur, etiam $2\equiv aa$(mod.$u$) \emph{i. e.} $2$ etiam ipsius $u$ residuum fore. Facile vero perspicitur, esse $u<t$, quare $t$ non est minimus numerus inductioni nostrae contrarius contra hyp. Unde manifesto sequitur id quod per inductionem inveneramus generaliter verum esse.

Combinando haec cum prop. art. 111 sequentia theoremata nanciscimur.

I. \emph{Numerorum omnium primorum formae $8n+3$, $+2$ erit non-residuum, $-2$ vero residuum.}

II. \emph{Numerorum omnium primorum formae $8n+5$ tum $+2$ tum $-2$ erunt non-residua.}

\Article{113}
Per similem inductionem ex tab. II inveniuntur numeri primi quorum residuum est $-2$ hi: $3,11,17,19,41,43,59,67,73,83,89,97$\HistoricalNote{*)}{Considerando scilicet $-2$ tamquam productum ex $+2$ et $-1$ V. art. 111.}. Inter quos quum nulli inveniantur formarum $8n+5$, $8n+7$, num etiam haec inductio theorematis generalis vim adipisci possit investigemus. Ostenditur simili modo ut in art. praec. quemvis numerum compositum formae $8n+5$ vel $8n+7$, factorem primum involvere formae $8n+5$ vel formae $8n+7$, ita ut, si inductio nostra generaliter vera, $-2$ nullius omnino numeri formae $8n+5$ vel $8n+7$ residuum esse possit. Si autem tales numeri darentur, ponatur omnium minimus $=t$, fiatque $-2=aa-tu$. Ubi si uti supra $a$ impar ipsoque $t$ minor accipitur, $u$ erit formae $8n+5$ vel $8n+7$, prout $t$ formae $8n+7$ vel $8n+5$. At ex eo quod $aa+2=tu$ atque $a<t$, quisquis facile derivare poterit, etiam
